% Section 1.6, from the summary table of lectures/L01/appendix/a_theory.md and from
% lectures/L01/README.md: the objectives, the evaluation questions and the next lesson.

\section{Sammanfattning}
\label{c1:sec:sammanfattning}

\begin{booktable}{@{}p{15em}L@{}}
  \thead{Begrepp} & \thead{Definition eller formel} \\
  \midrule
  Absolutbelopp & $\abs{a} \geq 0$, avståndet till noll \\
  Räkneordning & Parenteser, potenser, $\times$ och $\div$, $+$ och $-$ \\
  Bråkaddition & Gemensam nämnare \\
  Bråkmultiplikation & Täljare $\times$ täljare, nämnare $\times$ nämnare \\
  Procent & $p\,\% = \frac{p}{100}$ \\
  Parallellresistans & $\frac{1}{R_{\text{TOT}}} = \frac{1}{R_1} + \frac{1}{R_2}$ \\
  Parallellresistans, två motstånd & $R_{\text{p}} = \frac{R_1 \times R_2}{R_1 + R_2}$ \\
  Serieresistans & $R_{\text{s}} = R_1 + R_2$ \\
  Ohms lag & $U = R \times I$ \\
  Verkningsgrad & $\eta = \frac{P_{\text{ut}}}{P_{\text{in}}} \times 100\,\%$ \\
\end{booktable}

\needspace{6\baselineskip}
\subsection*{Det här ska du kunna}
Efter det här kapitlet ska du:
\begin{itemize}
  \item Känna till de olika talmängderna och deras egenskaper.
  \item Kunna tillämpa räkneordningen och de grundläggande räknesätten.
  \item Kunna räkna med bråk, decimaltal och procent.
  \item Kunna beräkna resistans, ström och spänning i enkla serie- och parallellkopplingar.
\end{itemize}

\testadigsjalv
\begin{enumerate}
  \item Till vilken talmängd hör $\sqrt{5}$? Motivera svaret kortfattat.
  \item Beräkna parallellresistansen för $R_1 = 4\,\Omega$ och $R_2 = 12\,\Omega$.
\end{enumerate}

\needspace{6\baselineskip}
\subsection*{Nästa kapitel}
\Kapref{2} handlar om algebra: algebraiska uttryck, förenkling och faktorisering.
